\documentclass[11pt]{article}

\usepackage[a4paper, margin=1in]{geometry}
\usepackage{amsmath}
\usepackage{booktabs}
\usepackage{hyperref}
\usepackage{url}
\usepackage{microtype}
\usepackage{enumitem}

\title{Supplement: MI-CLAIM checklist for the single-extractor
LLM endometriosis-Reddit study}
\date{\today}

\begin{document}
\maketitle

\noindent
This supplement walks through the Minimum Information about Clinical AI
Modeling (MI-CLAIM) checklist of Norgeot et al.~\cite{norgeot2020miclaim}
and addresses each item as it applies to our study. MI-CLAIM was written
with predictive clinical-AI systems in mind; the present work is an
LLM-based thematic-synthesis study of public patient-community text and
does not produce a clinical decision or a patient-level prediction. We
therefore include each MI-CLAIM section in full, mark items not
applicable (\emph{N/A}) where the construct does not map onto our
design, and justify each \emph{N/A} explicitly so reviewers can audit
our reasoning rather than treat the checklist as a presence/absence
exercise.

\section*{Part 1: Study Design}

\paragraph{Clinical problem and intended use.}
The study addresses an infodemiological question rather than a bedside
prediction: \emph{which themes in patient-authored r/Endo posts and
comments describe endometriosis care, and how reliable are the themes
from a completed single-extractor/subagent workflow?} The intended use
is to inform clinicians and researchers about under-discussed concerns
in online endometriosis communities, not to triage, diagnose, or treat
any individual patient. Outputs are descriptive themes, supporting
quotes, and per-theme reliability scores; no clinical action is
recommended on the basis of any model output.

\paragraph{Cohort and population.}
The ``cohort'' is a mentor-provided r/Endo archive containing
38{,}001 posts and 291{,}754 comments. The completed analysis used a
sample of 2{,}000 posts and 80{,}000 comments. Authors are
self-selected Reddit users; we make no claim that they are
representative of the global endometriosis population. The Limitations
section of the main manuscript discusses this explicitly.

\paragraph{Comparison group / baseline.}
The frozen baseline is the five-theme summary published in the 2024
GPT-4 abstract. We could not re-run GPT-4 because no OpenAI API
credential was available, so the comparison is necessarily a paper-vs-paper
comparison rather than a head-to-head re-run. This is documented as a
limitation. The Young~2015 systematic-review themes~\cite{young2015}
serve as a second, pre-LLM literature anchor against which any modern
extracted theme can be measured.

\paragraph{Outcome.}
The primary ``outcome'' is the per-theme composite reliability score in
$[0, 1]$ described in Section~2.4 of the main text. In the completed run
the composite averages the available normalized signals: quote
grounding, anchored validity against Young~2015, and LLM-as-judge
rubric scores from Sonnet~4.6 and Opus~4.7. There is no patient-level
clinical endpoint. Cross-model agreement, prompt-B stability,
high-temperature stability, bootstrap uncertainty, and subgroup
performance are not reported because those outputs were not produced.

\section*{Part 2: Data}

\paragraph{Data source.}
Posts and comments came from a mentor-provided r/Endo archive. The
completed run did not use a public-archive re-collection and does not
include other-subreddit subgroup claims.

\paragraph{Data preprocessing.}
The completed analysis samples 2{,}000 posts and 80{,}000 comments from
the 38{,}001-post, 291{,}754-comment archive. Text is normalized for
model input and quote checks, but no named-entity recognition,
re-identification, or personally identifying enrichment is performed.
No corpus-size sensitivity output is reported for the present run.

\paragraph{Data labelling.}
\emph{N/A.} The study has no class labels. The Young~2015 themes are
treated as an external anchor, not as supervised training labels. The
LLMs are used zero-shot.

\paragraph{Data partitioning.}
\emph{N/A in the supervised-learning sense.} No model is trained or
fine-tuned in this work, so there is no train/validation/test split.
The sampled corpus is partitioned only for inference: 74 overlapping
map chunks are processed independently before reduction. This is not an
independent validation split.

\paragraph{Missing data.}
Because the run uses a mentor-provided archive, we cannot quantify
archive coverage against live Reddit or recover content deleted before
the archive was provided. This is documented as a limitation. No
public-archive coverage or corpus-perturbation sensitivity result is
claimed.

\paragraph{Data ethics.}
The corpus is publicly accessible, pseudonymous Reddit content; the
study qualifies as non-human-subjects research under the public-data
category and was reviewed accordingly. No raw Reddit text is
redistributed; only SHA-256 manifests, prompts, model snapshots, and
aggregate results are released. All illustrative quotes in the
published manuscript are paraphrased to reduce re-identification risk,
in line with the open-science recommendations of
Norori et al.~\cite{norori2021addressing}.

\section*{Part 3: Model Development}

\paragraph{Model architecture.}
We use Claude Haiku 4.5~\cite{anthropic_claude_haiku} as the single
black-box thematic extractor. Claude Sonnet~4.6 performs the subagent
reduce step, and Claude Sonnet~4.6 plus Claude Opus~4.7 serve as LLM
judges. No Gemini extraction is included in the completed run.
Architectures are proprietary and not published; dated model snapshots
are recorded in the manuscript run index.

\paragraph{Software stack.}
The pipeline is implemented in Python. Quote grounding uses exact
substring checks and BM25 retrieval; reliability tables combine
grounding, Young~2015 anchor alignment, and Sonnet/Opus judge scores.
Exact code, pinned dependency versions, and prompts are released with
the repository.

\paragraph{Prompt and inference protocol.}
Theme extraction follows a map-reduce protocol. The sampled corpus is
packed into 74 chunks of approximately 100{,}000 tokens with a
2{,}000-token overlap. Claude Haiku 4.5 runs prompt~A at $T=0.0$ in the
map step via subagents and produces 722 chunk-level themes. A Claude
Sonnet~4.6 subagent reduce step merges those outputs into 20 final
themes. Prompt~B and high-temperature runs were not completed.

\paragraph{Hyperparameter selection.}
No model weights are tuned on the corpus. The completed run uses the
fixed operational settings above: prompt~A, $T=0.0$, 100{,}000-token
chunks, and 2{,}000-token overlap. Planned cross-model and
prompt/temperature perturbation settings are excluded from the reported
results because they did not complete.

\paragraph{Training compute and energy.}
\emph{N/A.} No model training was performed. Inference cost and token
counts per run are logged and reported as wall-clock and dollar cost in
the run index.

\section*{Part 4: Optimization and Validation}

\paragraph{Internal validation.}
Because we have no class labels, internal validation uses the completed
automated signals and explicitly marks planned-but-absent signals as
unavailable:
\begin{itemize}[leftmargin=1.2em]
  \item \textbf{Quote grounding (faithfulness).} Each supporting quote
    is checked against the raw corpus by exact substring match
    (whitespace-collapsed, case-insensitive) and BM25 retrieval. In the
    completed run, 21 of 60 supporting quotes are grounded (35\%).
  \item \textbf{Anchored validity.} The eight Young~2015 themes are
    encoded as a literature anchor. Nine of 20 final themes align with
    a Young anchor, covering seven of the eight anchors.
  \item \textbf{LLM-as-judge rubric.} Each theme is scored 1--5 on
    faithfulness, comprehensiveness, clinical relevance, specificity,
    and non-redundancy by Claude Sonnet~4.6 and Claude Opus~4.7. These
    models are judges, not independent extractors.
  \item \textbf{Chunk recurrence.} The reduce output records how often
    each final theme is supported across the 74 Haiku map chunks. This
    is descriptive support, not cross-model agreement.
  \item \textbf{Unavailable planned signals.} Gemini extraction,
    cross-model agreement, prompt-B agreement, high-temperature
    stability, bootstrap intervals, and year/subreddit subgroup outputs
    are not part of the completed results.
\end{itemize}
Per theme, grounding, judge, and anchor signals are normalized to
$[0, 1]$ and averaged into the composite reliability score. The
composite mean is 0.55 with range 0.28--0.87.

\paragraph{External validation.}
The Young~2015 systematic-review themes constitute the external,
pre-LLM, non-corpus reference. Convergent overlap with Young~2015 is
treated as supportive but not sufficient evidence; LLMs were trained on
text that includes the qualitative literature, so anchor overlap is
necessary but not sufficient for a claim of independent recovery. In
the completed run, nine of 20 final themes align with Young~2015 and
seven of the eight Young anchors are represented.

\paragraph{Sample size justification.}
The analysis sample contains 2{,}000 posts and 80{,}000 comments drawn
from a larger 38{,}001-post, 291{,}754-comment archive. This size was
chosen to keep the single-extractor/subagent run tractable while still
using a substantially comment-rich sample. No sample-size sensitivity
analysis is reported.

\section*{Part 5: Model Performance}

\paragraph{Performance metrics.}
The reportable performance object is the per-theme composite
reliability score and its completed components. The Sonnet reduce step
produced 20 final themes from 722 Haiku chunk-level themes. Quote
grounding is 35\% (21/60); Young alignment is 9/20 final themes and
7/8 anchors; the composite mean is 0.55 with range 0.28--0.87.

\paragraph{Confidence intervals / uncertainty.}
No confidence intervals are reported. Prompt-B runs, high-temperature
stability runs, and bootstrap resampling were planned but did not
produce outputs for the completed analysis. Theme-level uncertainty is
therefore summarized descriptively through grounding, anchor alignment,
judge scores, and chunk recurrence.

\paragraph{Calibration.}
\emph{N/A in the predictive-probability sense.} There is no
class-probability output to calibrate. The closest analogue is the
LLM-as-judge rubric, whose 1--5 scores are inherently subjective and
inherit the calibration limits of LLM judges; we discuss this in the
main-text Discussion.

\paragraph{Subgroup performance.}
Reddit metadata does not include patient demographics, and the completed
archive is r/Endo only. No year or subreddit subgroup outputs are
reported. Demographic subgroup analysis (age, race, geography,
comorbidity) is \emph{N/A} because the data do not support it.

\paragraph{Comparison to baseline.}
The 2024 GPT-4 five-theme summary is treated as a frozen baseline
for contextual comparison. It is not a live comparator because GPT-4
was not re-run. The completed results should be interpreted as a
single-extractor/subagent thematic synthesis, not as evidence of
cross-model recovery.

\section*{Part 6: Examination of Model}

\paragraph{Interpretability and explanation.}
Each extracted theme is accompanied by up to three verbatim quotes
and by a 2--3 sentence description; the grounding check tests whether
those quotes can be located in the corpus. In the completed run, 21 of
60 supporting quotes pass this lexical check, so quote fields remain a
priority for manual review before publication.

\paragraph{Failure mode analysis.}
We separately tabulate (a) themes with high LLM-as-judge scores but low
quote grounding, (b) themes with broad chunk recurrence but no Young
anchor, and (c) themes with low composite scores. Because the run has
one extractor, these analyses flag candidate single-extractor artefacts
rather than cross-model disagreements.

\paragraph{Bias and fairness audit.}
Reddit users are self-selected, English-speaking, and skew younger and
more digitally literate than the endometriosis population at large.
LLM judges are themselves a source of bias: each model's judgements of
another model's themes carry the calibration limits of LLM-as-judge
setups documented in the broader literature. We treat the reliability
score as a filtering aid, not a final-judgment substitute, and note the
absence of patient and clinician raters as the most important
limitation of the validation stack.

\paragraph{Deployment.}
\emph{N/A.} The pipeline is a one-off research artefact, not a
deployed clinical system. Re-execution against the same model
snapshots is supported by the released code, prompts, and run index;
re-execution against newer model snapshots is encouraged as a form of
longitudinal model audit.

\paragraph{Code and model availability.}
All code, prompts, dated model snapshots, the Young~2015 anchor
encoding, and aggregated reliability tables are released at the
repository URL listed in the Reproducibility statement of the main
text. Raw Reddit posts and comments are not redistributed; their
SHA-256 manifest is published so that collaborators with the same data
window can verify identical input.

\bibliographystyle{plain}
\bibliography{refs}

\end{document}
